\documentclass[10pt,a4paper]{amsart}
\usepackage{amsmath,amssymb,amsthm,graphicx}
\usepackage{verbatim}
\newenvironment{reference}{\comment}{\endcomment}
%\newenvironment{reference}{}{}
\newenvironment{slogan}{\comment}{\endcomment}
\newenvironment{history}{\comment}{\endcomment}
% For commutative diagrams we use Xy-pic
\usepackage[all]{xy}

% We use 2cell for 2-commutative diagrams.
\xyoption{2cell}
% Theorem environments.
%
\theoremstyle{plain}
\newtheorem{theorem}[subsection]{Theorem}
\newtheorem{proposition}[subsection]{Proposition}
\newtheorem{lemma}[subsection]{Lemma}

\theoremstyle{definition}
\newtheorem{definition}[subsection]{Definition}
\newtheorem{example}[subsection]{Example}
\newtheorem{exercise}[subsection]{Exercise}
\newtheorem{situation}[subsection]{Situation}

\theoremstyle{remark}
\newtheorem{remark}[subsection]{Remark}
\newtheorem{remarks}[subsection]{Remarks}

\numberwithin{equation}{subsection}

% Macros
%
\def\lim{\mathop{\mathrm{lim}}\nolimits}
\def\colim{\mathop{\mathrm{colim}}\nolimits}
\def\Spec{\mathop{\mathrm{Spec}}}
\def\Hom{\mathop{\mathrm{Hom}}\nolimits}
\def\Ext{\mathop{\mathrm{Ext}}\nolimits}
\def\SheafHom{\mathop{\mathcal{H}\!\mathit{om}}\nolimits}
\def\SheafExt{\mathop{\mathcal{E}\!\mathit{xt}}\nolimits}
\def\Sch{\mathit{Sch}}
\def\Mor{\mathop{\mathrm{Mor}}\nolimits}
\def\Ob{\mathop{\mathrm{Ob}}\nolimits}
\def\Sh{\mathop{\mathit{Sh}}\nolimits}
\def\NL{\mathop{N\!L}\nolimits}
\def\CH{\mathop{\mathrm{CH}}\nolimits}
\def\proetale{{pro\text{-}\acute{e}tale}}
\def\etale{{\acute{e}tale}}
\def\QCoh{\mathit{QCoh}}
\def\Ker{\mathop{\mathrm{Ker}}}
\def\Im{\mathop{\mathrm{Im}}}
\def\Coker{\mathop{\mathrm{Coker}}}
\def\Coim{\mathop{\mathrm{Coim}}}

% Boxtimes
%
\DeclareMathSymbol{\boxtimes}{\mathbin}{AMSa}{"02}

%
% Macros for moduli stacks/spaces
%
\def\QCohstack{\mathcal{QC}\!\mathit{oh}}
\def\Cohstack{\mathcal{C}\!\mathit{oh}}
\def\Spacesstack{\mathcal{S}\!\mathit{paces}}
\def\Quotfunctor{\mathrm{Quot}}
\def\Hilbfunctor{\mathrm{Hilb}}
\def\Curvesstack{\mathcal{C}\!\mathit{urves}}
\def\Polarizedstack{\mathcal{P}\!\mathit{olarized}}
\def\Complexesstack{\mathcal{C}\!\mathit{omplexes}}
% \Pic is the operator that assigns to X its picard group, usage \Pic(X)
% \Picardstack_{X/B} denotes the Picard stack of X over B
% \Picardfunctor_{X/B} denotes the Picard functor of X over B
\def\Pic{\mathop{\mathrm{Pic}}\nolimits}
\def\Picardstack{\mathcal{P}\!\mathit{ic}}
\def\Picardfunctor{\mathrm{Pic}}
\def\Deformationcategory{\mathcal{D}\!\mathit{ef}}
\usepackage{cleveref}
\crefname{originallemma}{original lemma}{original lemmas}
\Crefname{originallemma}{Original lemma}{Original lemmas}
\makeatletter
\newlabel{lemma-closed-in-quasi-compact}{{lemma-closed-in-quasi-compact}{source}{}{sourcepointer-1}{}}
\newlabel{lemma-closed-in-quasi-compact@cref}{{[originallemma][0][]lemma-closed-in-quasi-compact}{[1][0][]source}{}{}{}}
\makeatother
\crefname{originallemma}{original lemma}{original lemmas}
\Crefname{originallemma}{Original lemma}{Original lemmas}
\makeatletter
\newlabel{lemma-connected-components}{{lemma-connected-components}{source}{}{sourcepointer-2}{}}
\newlabel{lemma-connected-components@cref}{{[originallemma][0][]lemma-connected-components}{[1][0][]source}{}{}{}}
\makeatother
\crefname{originallemma}{original lemma}{original lemmas}
\Crefname{originallemma}{Original lemma}{Original lemmas}
\makeatletter
\newlabel{lemma-image-quasi-compact}{{lemma-image-quasi-compact}{source}{}{sourcepointer-3}{}}
\newlabel{lemma-image-quasi-compact@cref}{{[originallemma][0][]lemma-image-quasi-compact}{[1][0][]source}{}{}{}}
\makeatother
\crefname{originallemma}{original lemma}{original lemmas}
\Crefname{originallemma}{Original lemma}{Original lemmas}
\makeatletter
\newlabel{lemma-intersection-closed-in-quasi-compact}{{lemma-intersection-closed-in-quasi-compact}{source}{}{sourcepointer-4}{}}
\newlabel{lemma-intersection-closed-in-quasi-compact@cref}{{[originallemma][0][]lemma-intersection-closed-in-quasi-compact}{[1][0][]source}{}{}{}}
\makeatother
\crefname{originallemma}{original lemma}{original lemmas}
\Crefname{originallemma}{Original lemma}{Original lemmas}
\makeatletter
\newlabel{lemma-profinite}{{lemma-profinite}{source}{}{sourcepointer-5}{}}
\newlabel{lemma-profinite@cref}{{[originallemma][0][]lemma-profinite}{[1][0][]source}{}{}{}}
\makeatother
\crefname{originallemma}{original lemma}{original lemmas}
\Crefname{originallemma}{Original lemma}{Original lemmas}
\makeatletter
\newlabel{lemma-space-connected-components}{{lemma-space-connected-components}{source}{}{sourcepointer-6}{}}
\newlabel{lemma-space-connected-components@cref}{{[originallemma][0][]lemma-space-connected-components}{[1][0][]source}{}{}{}}
\makeatother
\title{The space of connected components of a quasi-compact space whose components are intersections of clopens is profinite}
\author{The Stacks Project Authors}
\date{stacks-project:topology.tex:local-snapshot-20261012; GNU FDL 1.2 or later}
\begin{document}
\maketitle
\noindent\textit{Editorial note.} All delivered mathematical text is the author's own, copied from the pinned Stacks Project chapter source sources/stacks-topology/pinned-original-source.tex; no statement or proof step was rewritten, reordered or removed.\par\medskip
\noindent\textit{Editorial note.} Preamble provenance: the Stacks Project keeps one single shared preamble for every chapter; this card takes it directly from the staged sources/stacks-topology/preamble.tex as three verbatim spans (lines 8-12, 14-18 and 154-217), so no derived file was written and nothing was added or redefined. The wrapper class and the base packages are supplied by the assembler.\par\medskip
\noindent\textit{Editorial note.} Bibliography provenance: neither the counted statement nor its proof contains a citation (this lane rescans the whole span before writing the card), so the card carries no bibliography and no entry was reconstructed.\par\medskip
\noindent\textit{Editorial note.} Selection provenance: this card was opened under the file-based control contract (CONTROL.md), with the mechanical pre-card gate run on its statement and proof interval before the card was written.\par\medskip
\noindent\textit{Allowed standard prerequisites (editorial).}\begin{itemize}
\item The characterization of profinite spaces as Hausdorff quasi-compact totally disconnected. Invoked at the start of the counted proof (source line 4457). Reduces the assertion to the three properties.
\item The statement that the image of a quasi-compact space is quasi-compact, and that the space of components is totally disconnected by construction. Invoked inside the counted proof (source lines 4458-4461). Gives quasi-compactness and total disconnectedness.
\item The statement that a component is the intersection of the clopens containing it, together with the finite intersection criterion. Invoked inside the counted proof (source lines 4463-4468). Produces a clopen separating two components.
\end{itemize}\medskip
\ifdefined\hypertarget\hypertarget{sourcepointer-1}{}\fi
\noindent\textit{Source reference.} Original lemma lemma-closed-in-quasi-compact: sources/stacks-topology/pinned-original-source.tex, source lines 1879--1879. Navigation pointer only: the counted text refers to the Stacks item lemma-closed-in-quasi-compact, which is an existing labeled result of this same chapter and is not part of this card. The pointer carries the target's own label as its number and its definition line as its span; its content is not reproduced here.\par\medskip
\ifdefined\hypertarget\hypertarget{sourcepointer-2}{}\fi
\noindent\textit{Source reference.} Original lemma lemma-connected-components: sources/stacks-topology/pinned-original-source.tex, source lines 875--875. Navigation pointer only: the counted text refers to the Stacks item lemma-connected-components, which is an existing labeled result of this same chapter and is not part of this card. The pointer carries the target's own label as its number and its definition line as its span; its content is not reproduced here.\par\medskip
\ifdefined\hypertarget\hypertarget{sourcepointer-3}{}\fi
\noindent\textit{Source reference.} Original lemma lemma-image-quasi-compact: sources/stacks-topology/pinned-original-source.tex, source lines 1964--1964. Navigation pointer only: the counted text refers to the Stacks item lemma-image-quasi-compact, which is an existing labeled result of this same chapter and is not part of this card. The pointer carries the target's own label as its number and its definition line as its span; its content is not reproduced here.\par\medskip
\ifdefined\hypertarget\hypertarget{sourcepointer-4}{}\fi
\noindent\textit{Source reference.} Original lemma lemma-intersection-closed-in-quasi-compact: sources/stacks-topology/pinned-original-source.tex, source lines 1941--1941. Navigation pointer only: the counted text refers to the Stacks item lemma-intersection-closed-in-quasi-compact, which is an existing labeled result of this same chapter and is not part of this card. The pointer carries the target's own label as its number and its definition line as its span; its content is not reproduced here.\par\medskip
\ifdefined\hypertarget\hypertarget{sourcepointer-5}{}\fi
\noindent\textit{Source reference.} Original lemma lemma-profinite: sources/stacks-topology/pinned-original-source.tex, source lines 4351--4351. Navigation pointer only: the counted text refers to the Stacks item lemma-profinite, which is an existing labeled result of this same chapter and is not part of this card. The pointer carries the target's own label as its number and its definition line as its span; its content is not reproduced here.\par\medskip
\ifdefined\hypertarget\hypertarget{sourcepointer-6}{}\fi
\noindent\textit{Source reference.} Original lemma lemma-space-connected-components: sources/stacks-topology/pinned-original-source.tex, source lines 1014--1014. Navigation pointer only: the counted text refers to the Stacks item lemma-space-connected-components, which is an existing labeled result of this same chapter and is not part of this card. The pointer carries the target's own label as its number and its definition line as its span; its content is not reproduced here.\par\medskip
\begin{lemma}
\label{lemma-pi0-profinite}
Let $X$ be a topological space. If $X$ is quasi-compact
and every connected component of $X$ is the intersection
of the open and closed subsets containing it, then $\pi_0(X)$
is a profinite space.
\end{lemma}

\begin{proof}
We will use Lemma \ref{lemma-profinite} to prove this.
Since $\pi_0(X)$ is the image of a quasi-compact space it is
quasi-compact (Lemma \ref{lemma-image-quasi-compact}).
It is totally disconnected by construction
(Lemma \ref{lemma-space-connected-components}).
Let $C, D \subset X$ be distinct connected components of $X$.
Write $C = \bigcap U_\alpha$ as the intersection of the open and
closed subsets of $X$ containing $C$. Any finite intersection
of $U_\alpha$'s is another. Since $\bigcap U_\alpha \cap D = \emptyset$
we conclude that $U_\alpha \cap D = \emptyset$ for some $\alpha$
(use Lemmas \ref{lemma-connected-components},
\ref{lemma-closed-in-quasi-compact} and
\ref{lemma-intersection-closed-in-quasi-compact})
Since $U_\alpha$ is open and closed, it is the union of the
connected components it contains, i.e., $U_\alpha$ is the inverse
image of some open and closed subset $V_\alpha \subset \pi_0(X)$.
This proves that the points corresponding to $C$ and $D$
are contained in disjoint open subsets, i.e., $\pi_0(X)$ is
Hausdorff.
\end{proof}

\end{document}
